\subsection{Exactness and Diagrams}

\begin{theorem} \label{thm:first-isomorphism-theorem-vector-spaces} \hlt{First Isomorphism Theorem}\\
Let $T\in\LL(\mathcal{V},\mathcal{W})$. Then $\widehat{T}:\mathcal{V}/\ker(T)\to\operatorname{range}(T)$; $\widehat{T}(\mathbf{v}+\ker(T))=T\mathbf{v}$ is a well-defined linear isomorphism.
\begin{proof}
If $\mathbf{v}+\ker(T)=\mathbf{w}+\ker(T)$, then $\mathbf{v}-\mathbf{w}\in\ker(T)$, so $T\mathbf{v}=T\mathbf{w}$; hence $\widehat{T}$ is well-defined.\\
Linearity follows from the quotient operations, $\widehat{T}(a(\mathbf{v}+\ker(T))+(\mathbf{w}+\ker(T)))=aT\mathbf{v}+T\mathbf{w}$.\\
Also, $\widehat{T}(\mathbf{v}+\ker(T))=\mathbf{0} \iff T\mathbf{v}=\mathbf{0} \iff \mathbf{v}\in\ker(T) \iff \mathbf{v}+\ker(T)=\ker(T)$, so $\widehat{T}$ is injective.\\
If $\mathbf{y}\in\operatorname{range}(T)$, then $\mathbf{y}=T\mathbf{v}=\widehat{T}(\mathbf{v}+\ker(T))$ for some $\mathbf{v}\in\mathcal{V}$, so $\widehat{T}$ is surjective. Thus $\widehat{T}$ is linear isomorphism.\\
\hlt{Structural Commentary}: The quotient by the kernel is the part of the domain actually seen by $T$.\\
\hlt{Key Insight}: Image is domain modulo kernel.
\end{proof}
\end{theorem}

\begin{theorem} \label{thm:second-isomorphism-vector-spaces} \hlt{Second Isomorphism Theorem}\\
Let $\mathcal{S}$ and $\mathcal{T}$ be subspaces of $\mathcal{V}$. Then $(\mathcal{S}+\mathcal{T})/\mathcal{T}\cong \mathcal{S}/(\mathcal{S}\cap\mathcal{T})$.
\begin{proof}
Define $L:\mathcal{S}\to(\mathcal{S}+\mathcal{T})/\mathcal{T}$ by $L\mathbf{s}=\mathbf{s}+\mathcal{T}$. Then $L$ is linear by the quotient operations.\\
Moreover, $\ker(L)=\{\mathbf{s}\in\mathcal{S}:\mathbf{s}+\mathcal{T}=\mathcal{T}\} =\mathcal{S}\cap\mathcal{T}$.
If $(\mathbf{s}+\mathbf{t})+\mathcal{T}\in(\mathcal{S}+\mathcal{T})/\mathcal{T}$, with $\mathbf{s}\in\mathcal{S}$ and $\mathbf{t}\in\mathcal{T}$, then $(\mathbf{s}+\mathbf{t})+\mathcal{T}=\mathbf{s}+\mathcal{T}=L\mathbf{s}$, so $L$ is surjective. Theorem \ref{thm:first-isomorphism-theorem-vector-spaces} gives the stated isomorphism.\\
\hlt{Structural Commentary}: Modding out by $\mathcal{T}$ leaves only the part of $\mathcal{S}$ not already inside $\mathcal{T}$.\\
\hlt{Key Insight}: In a quotient by $\mathcal{T}$, adding $\mathcal{T}$ to $\mathcal{S}$ costs nothing.
\end{proof}
\end{theorem}

\begin{theorem} \label{thm:third-isomorphism-vector-spaces} \hlt{Third Isomorphism Theorem}\\
If $\mathcal{U}\subseteq\mathcal{T}\subseteq\mathcal{V}$ are subspaces, then $(\mathcal{V}/\mathcal{U})/(\mathcal{T}/\mathcal{U})\cong \mathcal{V}/\mathcal{T}$.
\begin{proof}
Define $L:\mathcal{V}/\mathcal{U}\to\mathcal{V}/\mathcal{T}$ by $L(\mathbf{v}+\mathcal{U})=\mathbf{v}+\mathcal{T}$.\\
If $\mathbf{v}+\mathcal{U}=\mathbf{w}+\mathcal{U}$, then $\mathbf{v}-\mathbf{w}\in\mathcal{U}\subseteq\mathcal{T}$, so $L$ is well-defined.\\
It is linear by the quotient operations, and it is surjective since $L(\mathbf{v}+\mathcal{U})=\mathbf{v}+\mathcal{T}$ for every coset $\mathbf{v}+\mathcal{T}$.\\
Moreover, $\ker(L)=\{\mathbf{v}+\mathcal{U}: \mathbf{v}+\mathcal{T}=\mathcal{T}\} =\{\mathbf{v}+\mathcal{U}: \mathbf{v}\in\mathcal{T}\} =\mathcal{T}/\mathcal{U}$. Theorem \ref{thm:first-isomorphism-theorem-vector-spaces} gives $(\mathcal{V}/\mathcal{U})/(\mathcal{T}/\mathcal{U})\cong\mathcal{V}/\mathcal{T}$.\\
\hlt{Structural Commentary}: Collapsing $\mathcal{U}$ and then collapsing $\mathcal{T}/\mathcal{U}$ is the same as collapsing $\mathcal{T}$ at once.\\
\hlt{Key Insight}: Quotients of quotients cancel only when the nesting is explicit.
\end{proof}
\end{theorem}
